\documentclass[11pt]{amsart}
\usepackage{amsmath,amssymb,amsthm}
\usepackage[margin=1.2in]{geometry}
\usepackage{hyperref}

\newtheorem{theorem}{Theorem}[section]
\newtheorem{lemma}[theorem]{Lemma}
\newtheorem{proposition}[theorem]{Proposition}
\newtheorem{corollary}[theorem]{Corollary}
\newtheorem{conjecture}[theorem]{Conjecture}
\theoremstyle{definition}
\newtheorem{definition}[theorem]{Definition}
\newtheorem{remark}[theorem]{Remark}
\newtheorem{example}[theorem]{Example}

\newcommand{\PF}{\mathcal{P}}
\newcommand{\ZZ}{\mathbb{Z}}
\newcommand{\bott}[1]{\mathrm{bot}_{#1}}
\newcommand{\topp}[1]{\mathrm{top}_{#1}}

\title[The canonical box partition of $G$-parking functions]{A canonical box
partition of the set of $G$-parking functions,\\ with an application to the
spanning trees of the hypercube}
\author{Working note --- QnParkingFunctions project}
\date{Draft of August 19, 2026}

\begin{document}
\maketitle

\begin{abstract}
For a connected multigraph $G$ with root $q$ and a total order on the
vertices, we show that the fibers of the map sending each $G$-parking
function to its lexicographically minimal Dhar burn order are coordinate
boxes. The boxes partition the set of all $G$-parking functions; their tops
are exactly the maximal $G$-parking functions. This yields sign-free formulas
for $T_G(1,1)$ and $T_G(1,y)$ as sums of products of $q$-integers over
acyclic orientations with unique sink, gives an explicit and canonical
construction of an $M$-partition of the parking-function multicomplex whose
existence follows from Merino's inductive proof of Chari's conjecture for
cographic matroids, and, for the hypercube $Q_n$, contains the canonical box
of Benson--Chakrabarty--Tetali as its largest cell. We also prove a
support-product theorem for the spanning trees of $G\times K_2$ whose
iteration telescopes to Stanley's formula for the number of spanning trees
of $Q_n$, give an explicit bijection realizing its full-support slice, and
state precisely the bijective step that remains open.
\end{abstract}

\section{Setup}

Let $G=(V,E)$ be a connected multigraph with $n$ vertices, a distinguished
root $q$, and a fixed total order on $V$ with $q$ smallest; we write
$V=\{v : 0 \le v \le n-1\}$ with $q=0$. For $u,v \in V$ let $a(u,v)$ denote
the number of edges between $u$ and $v$.

A function $f\colon V\setminus\{q\} \to \ZZ_{\ge 0}$ is a
\emph{$G$-parking function} (equivalently, a \emph{superstable
configuration}) if for every nonempty $A \subseteq V\setminus\{q\}$ there is
$v\in A$ with $f(v) < \sum_{u \notin A} a(v,u)$. We write $\PF(G,q)$ for the
set of $G$-parking functions and recall Dhar's criterion
\cite{dhar,bct}: $f \in \PF(G,q)$ if and only if the following process burns
every vertex. Initially only $q$ is burnt; a vertex $v$ may burn whenever
\[
  f(v) \;<\; w(v) \;:=\; \sum_{u \text{ burnt}} a(v,u),
\]
and the process is abelian: if some order burns everything, every maximal
sequence of legal burns does.

By downward closure (\cite[Prop.~2.2]{bct}) $\PF(G,q)$ is an order ideal of
$\ZZ_{\ge 0}^{V \setminus \{q\}}$ under the coordinatewise order; all its
maximal elements have the same total degree $g := |E|-|V|+1$
(\cite[Cor.~3.2]{bct}), they are in bijection with the acyclic orientations
of $G$ with unique source at $q$ (\cite[Thm.~3.1]{bct}), and their number is
the Tutte evaluation $T_G(1,0)$ (\cite[Thm.~4.1]{bct}).

\begin{definition}\label{def:canon}
For $f \in \PF(G,q)$, the \emph{canonical burn order} $\sigma(f)$ is the
sequence obtained by always burning the \emph{smallest} burnable vertex.
By the abelian property $\sigma(f)$ is a well-defined ordering
$(\sigma_0=q,\sigma_1,\dots,\sigma_{n-1})$ of $V$.
\end{definition}

For an arbitrary ordering $\sigma$ of $V$ starting at $q$ and $0\le k<n$
write $w_k(v) = \sum_{j<k} a(v,\sigma_j)$ for the burnt-neighbor weight of
$v$ before step $k$. Define, for $v=\sigma_k$,
\[
  \topp{\sigma}(v) = w_k(v) - 1,
  \qquad
  \bott{\sigma}(v) = \max\bigl(\{0\} \cup \{\, w_j(v) : j<k,\ \sigma_j > v \,\}\bigr).
\]
Since $w_\bullet(v)$ is nondecreasing, $\bott{\sigma}(v)$ is the
burnt-neighbor weight of $v$ at the \emph{last} step that burned a larger
vertex while $v$ was unburnt (its ``skip watermark''), and $0$ if there was
no such step.

\section{The fiber theorem}

\begin{theorem}\label{thm:fiber}
Let $f \in \PF(G,q)$ and $\sigma = \sigma(f)$. Then
\[
  \{\, g \in \PF(G,q) \;:\; \sigma(g) = \sigma \,\}
  \;=\;
  \{\, g \;:\; \bott{\sigma} \le g \le \topp{\sigma} \ \text{coordinatewise} \,\},
\]
and moreover every integer vector in the right-hand box is a $G$-parking
function.
\end{theorem}

\begin{proof}
($\subseteq$) Let $g\in\PF(G,q)$ with $\sigma(g)=\sigma$ and let
$v=\sigma_k$. Since $v$ burns at step $k$ we have $g(v) < w_k(v)$, i.e.\
$g(v)\le\topp{\sigma}(v)$. If $j<k$ and $\sigma_j > v$, then at step $j$ the
canonical process preferred the larger vertex $\sigma_j$, so $v$ was not
burnable: $g(v) \ge w_j(v)$. Taking the maximum over such $j$ gives
$g(v)\ge\bott{\sigma}(v)$.

($\supseteq$) Let $g$ be any integer vector with
$\bott{\sigma}\le g\le\topp{\sigma}$. We show by induction on $k$ that the
canonical process for $g$ burns $\sigma_1,\dots,\sigma_{k}$ in this order.
Assume it has burnt exactly $\{\sigma_0,\dots,\sigma_{k-1}\}$, so each
unburnt $v$ has burnt-neighbor weight $w_k(v)$. First, $v=\sigma_k$ is
burnable: $g(v) \le \topp{\sigma}(v) = w_k(v)-1 < w_k(v)$. Second, let
$u < \sigma_k$ be unburnt; then $u = \sigma_m$ for some $m > k$, so $u$ is
skipped at step $k$ in $\sigma$, whence
$\bott{\sigma}(u) \ge w_k(u)$ and therefore $g(u) \ge w_k(u)$: $u$ is not
burnable. Hence the smallest burnable vertex is exactly $\sigma_k$. By
induction the process burns all of $V$, so $g \in \PF(G,q)$ and
$\sigma(g)=\sigma$.
\end{proof}

\begin{corollary}\label{cor:partition}
The boxes $B_\sigma := [\bott{\sigma},\topp{\sigma}]$, over the canonical burn
orders $\sigma$ of parking functions, partition $\PF(G,q)$. The map
$\sigma \mapsto \topp{\sigma}$ is a bijection onto the set of maximal
$G$-parking functions; in particular the number of boxes is $T_G(1,0)$.
\end{corollary}

\begin{proof}
The fibers of $f \mapsto \sigma(f)$ partition $\PF(G,q)$, and by
Theorem~\ref{thm:fiber} each fiber is a box. Every edge of $G$ contributes
its multiplicity to $w_k(\sigma_k)$ exactly once, at the step burning its
later endpoint, so
$\sum_v \topp{\sigma}(v) = |E| - (n-1) = g$; hence each top is a maximal
parking function \cite[Prop.~2.3, Cor.~3.2]{bct}. The top lies in its own
box, so $\sigma(\topp{\sigma})=\sigma$ and $\sigma\mapsto\topp{\sigma}$ is
injective. If $m$ is maximal then $m \in B_{\sigma(m)}$, so
$m \le \topp{\sigma(m)}$; both have degree $g$, forcing
$m=\topp{\sigma(m)}$. The count is \cite[Thm.~4.1]{bct}.
\end{proof}

\begin{corollary}[sign-free Tutte formulas]\label{cor:tutte}
Writing $[k]_y = 1+y+\dots+y^{k-1}$ and $b_m := \bott{\sigma(m)}$ for a
maximal parking function $m$,
\[
  T_G(1,y) \;=\; \sum_{m \text{ maximal}} \ \prod_{v \ne q}
      \bigl[\, m(v) - b_m(v) + 1 \,\bigr]_y ,
  \qquad
  T_G(1,1) \;=\; \sum_{m} \prod_{v\ne q} \bigl(m(v)-b_m(v)+1\bigr).
\]
\end{corollary}

\begin{proof}
By Merino's theorem \cite{merino97}, $T_G(1,y)=\sum_{f\in\PF(G,q)}
y^{\,g-\deg f}$. Group the sum by the boxes of
Corollary~\ref{cor:partition}: within $B_\sigma$, the complement
$\topp{\sigma}(v)-f(v)$ ranges independently over
$\{0,\dots,\topp{\sigma}(v)-\bott{\sigma}(v)\}$.
\end{proof}

Via \cite[Thm.~3.1]{bct} the index set is the set of acyclic orientations
with unique source at $q$, and $m(v)+1$ is the weighted in-degree of $v$;
the watermark $b_m$ is a statistic on such orientations that appears to be
new. Corollary~\ref{cor:tutte} replaces the $2^k$-term inclusion--exclusion
of \cite[Cor.~3.3]{bct} with a sum of $T_G(1,0)$ positive terms.

\begin{proposition}[the identity cell]\label{prop:identity}
Suppose the vertex labeling is such that every prefix
$\{0,1,\dots,k\}$ induces a connected subgraph. Then the identity ordering
$\iota=(0,1,\dots,n-1)$ is a canonical burn order, $\bott{\iota}\equiv 0$, and
$\topp{\iota}(v)=\bigl(\sum_{u<v}a(v,u)\bigr)-1$. For the hypercube $Q_n$
with its binary labeling, $\topp{\iota}(v)=\mathrm{wt}(v)-1$ is the
canonical parking function of \cite[Thm.~5.2]{bct} and the cell $B_\iota$ is
its entire lower box, of cardinality $\prod_{v\neq 0}\mathrm{wt}(v)
= \prod_{k=2}^n k^{\binom nk}$.
\end{proposition}

\begin{proof}
Burning in increasing label order never skips a smaller unburnt vertex, so
no watermark is ever raised and each step is legal because $v$ has a burnt
(smaller) neighbor by the prefix-connectivity assumption. In $Q_n$ the
smaller neighbors of $v$ are exactly the $\mathrm{wt}(v)$ vertices obtained
by clearing one bit.
\end{proof}

\begin{conjecture}[$M$-shelling]\label{conj:shelling}
In the lexicographic order of the canonical burn orders, every initial union
of boxes of Corollary~\ref{cor:partition} is an order ideal; that is, the
partition together with this order is an $M$-shelling in the sense of Chari
\cite{chari}.
\end{conjecture}

Conjecture~\ref{conj:shelling} has been verified exhaustively for $Q_2$,
$Q_3$, a $4$-vertex multigraph, and $25$ pseudorandom connected multigraphs
on up to $6$ vertices. The \emph{existence} of an $M$-shelling of this
multicomplex is a theorem of Merino \cite[Thm.~5.8]{merino01}, proved by
induction on parallel classes; the construction above is explicit and
appears to be the first closed-form one (a sweep of the citation trails of
\cite{merino01,chari}, of the Gr\"obner-theoretic parking-function
literature, and of the sandpile literature through 2026 found no interval
partition of $\PF(G,q)$ for general $G$). Two pieces of prior art deserve
emphasis: for complete graphs, the fibers of the outcome map of classical
parking functions are known to be boxes with a product formula
\cite[Prop.~3.3]{colmenarejo} --- Theorem~\ref{thm:fiber} specializes to
this; and the greedy-burning engine underlying $\sigma(f)$ is the same one
Cori and Le~Borgne used for their activity-preserving bijection to spanning
trees \cite{corilb}, of which the present partition can be regarded as the
fiber-level strengthening. On the algorithmic side, the search tree of
Theorem~\ref{thm:fiber} enumerates the maximal parking functions of $G$
(equivalently, by \cite[Thm.~3.1]{bct}, the acyclic orientations with unique
fixed source) one per leaf with no duplicate suppression; we are not aware
of a published enumeration algorithm stated for maximal parking functions.

\section{The support-product conjecture for $Q_n$}

Write $Q_n = Q_{n-1} \times K_2$ and fix the direction $d=n$ (all directions
are equivalent). A spanning tree $T$ of $Q_n$ contains, for each
$u \in V(Q_{n-1})$, at most one \emph{vertical} edge
$\{(u,0),(u,1)\}$; let $S(T) \subseteq V(Q_{n-1})$ be the set of $u$ for
which it does. For $S \subseteq V(Q_{n-1})$ let $F_{Q_{n-1}}(S)$ denote the
number of spanning forests of $Q_{n-1}$ with exactly $|S|$ components, one
containing each vertex of $S$ (equivalently, by the all-minors Matrix-Tree
theorem, the determinant of the Laplacian of $Q_{n-1}$ with the rows and
columns of $S$ deleted).

\begin{theorem}[support product]\label{conj:support}
Let $G$ be any connected graph with $M$ vertices and $c(G)$ spanning trees.
For every nonempty $S \subseteq V(G)$, the number of spanning trees of
$G \times K_2$ with vertical support exactly $S$ is
$\tfrac{1}{2}\,c(G)\, 2^{|S|} F_G(S)$. In particular, for $G = Q_{n-1}$ the
number of pairs $(T,r)$ with $T$ a spanning tree of $Q_n$, $r$ a root, and
$S(T)=S$ equals
\[
  N(S) \;=\; 2^{\,|S|}\cdot F_{Q_{n-1}}(S)\cdot 2^{\,n-1} c(Q_{n-1}).
\]
\end{theorem}

\begin{proof}
Give the vertical edge at $u$ the weight $z_u$ and every horizontal edge
weight $1$, and let $t(z) = \sum_T \prod_{u \in S(T)} z_u$ be the weighted
spanning-tree sum of $G\times K_2$. Order the vertices as two level-blocks;
the weighted Laplacian is
$\bigl(\begin{smallmatrix} L_G + Z & -Z \\ -Z & L_G + Z\end{smallmatrix}\bigr)$
with $Z = \operatorname{diag}(z_u)$. Conjugating by the orthogonal
level-swap basis $\frac{1}{\sqrt 2}\bigl(\begin{smallmatrix} I & I \\ I &
-I \end{smallmatrix}\bigr)$ block-diagonalizes it as $L_G \oplus (L_G+2Z)$.
The product of the nonzero Laplacian eigenvalues of a weighted graph on $N$
vertices equals $N$ times its tree sum; the zero eigenvalue lives in the
$L_G$ block, so with $N = 2M$,
\[
  t(z) \;=\; \frac{1}{2M}\Bigl(\prod_{\lambda \neq 0}\lambda(L_G)\Bigr)
             \det(L_G + 2Z)
        \;=\; \frac{c(G)}{2}\,\det(L_G + 2Z).
\]
Expanding the determinant by principal minors in the diagonal perturbation
$2Z$ and applying the all-minors Matrix-Tree theorem,
$\det(L_G+2Z) = \sum_{R} 2^{|R|} F_G(R) \prod_{u\in R} z_u$. Both sides of
$t(z)$ are multilinear in the $z_u$, and comparing the coefficient of
$\prod_{u\in S} z_u$ gives the claim; multiplying by the $2M = 2^n$ root
choices gives the hypercube form.
\end{proof}

\begin{proposition}[full-support slice, bijectively]\label{prop:fullsupport}
For $S = V(G)$ the theorem has an explicit bijection: a spanning tree of
$G\times K_2$ containing every vertical edge is exactly a choice of a
spanning tree of $G$ (its image under contracting all vertical edges)
together with a free level in $\{0,1\}$ for each of its $M-1$ edges; hence
there are $c(G)\,2^{M-1} = \tfrac12 c(G) 2^{|S|} F_G(S)$ of them, as
$F_G(V(G))=1$.
\end{proposition}

\begin{remark}[partial doubles]\label{rem:partialdouble}
For general $S$, contracting the vertical edges of a tree with support $S$
gives a bijection with the spanning trees of the \emph{partial double}
$G^{(S)}$: the graph obtained from $G$ by splitting every vertex $u \notin
S$ into two copies $u_0, u_1$ (an edge $(u,v)$ becomes two parallel edges if
$u,v\in S$; the pair $(u,v_0),(u,v_1)$ if only $u \in S$; and
$(u_0,v_0),(u_1,v_1)$ if $u,v\notin S$). Theorem~\ref{conj:support} is thus
equivalent to the identity $c(G^{(S)}) = \tfrac12 c(G) 2^{|S|} F_G(S)$, and
the remaining \emph{bijective} problem is to realize this identity by an
explicit map --- for instance by splitting one vertex at a time.
\end{remark}

\begin{corollary}\label{prop:telescope}
Stanley's formula
$c(Q_n) = 2^{2^n-n-1}\prod_{k=1}^n k^{\binom nk}$ follows by iterating
Theorem~\ref{conj:support} over the $n$ directions.
\end{corollary}

\begin{proof}
Summing over $S$ and using
$\sum_{S} z^{|S|} F_G(S) = \det(zI + L_G)$ (all-minors Matrix-Tree; the
$S=\varnothing$ term vanishes) at $z=2$,
\[
  2^n c(Q_n) \;=\; \sum_{S\neq\varnothing} N(S)
  \;=\; 2^{n-1}c(Q_{n-1}) \cdot \det\bigl(2I + L_{Q_{n-1}}\bigr)
  \;=\; 2^{n-1}c(Q_{n-1}) \prod_{k=0}^{n-1} (2k+2)^{\binom{n-1}{k}} .
\]
Writing $R_n := 2^n c(Q_n)$, this is
$R_n = R_{n-1} \prod_{j=1}^{n} (2j)^{\binom{n-1}{j-1}}$, and by induction
with Pascal's rule $\binom{n-1}{j}+\binom{n-1}{j-1}=\binom nj$ one obtains
$R_n = \prod_{j=1}^n (2j)^{\binom nj}$, which is Stanley's formula.
\end{proof}

Theorem~\ref{conj:support} was found experimentally and confirmed
\emph{exactly, class by class}, for $n=2$, for $n=3$ (all $15$ supports),
and for $n=4$ (all $255$ supports, from the complete set of
$679\,477\,248$ spanning-tree--root pairs) before the proof above was
identified; the identity is elementary by weighted Matrix-Tree arguments,
and closely related product-graph computations appear in
\cite{bernardi} and in Martin--Reiner, so we regard the \emph{statement} as
possibly folklore --- the content for this project is its bijective
realization. Two supporting facts, verified
computationally: (i) conditioned on $S$, the spins of the vertical edges
(the side of each vertical edge nearer the root) are \emph{exactly} jointly
uniform over $\{0,1\}^S$ for the ensemble of (tree, root) pairs --- this is
a transported and marginalized form of Bernardi's independence theorem
\cite[Thm.~1]{bernardi}, and accounts for the factor $2^{|S|}$; (ii) the
uniformity fails for trees with a fixed root, so any bijective realization
must work with rooted-anywhere trees and only afterwards quotient by the
$2^n$ root choices.

A bijective proof of Theorem~\ref{conj:support} would amount to a
bijection
\[
  \{(T,r) : S(T)=S\} \;\longleftrightarrow\;
  \{0,1\}^S \times \mathcal F_{Q_{n-1}}(S) \times
  \{\text{rooted spanning trees of } Q_{n-1}\},
\]
and, by Proposition~\ref{prop:telescope}, its iteration over $n$ would
constitute a fully bijective proof of Stanley's formula --- the problem
stated as open in \cite[Example 5.6.10]{ec2} and, for the bijective form,
still open after \cite{bernardi} (cf.\ the Notes to Chapter~5 of the second
edition of \cite{ec2}).

A structural reformulation: deleting the vertical edges of $T$ leaves
forests $H_0, H_1$ in the two levels, and $T$ is a tree if and only if the
bipartite ``component graph'' on the components of $H_0$ and $H_1$, with one
edge per element of $S$, is a tree. Conjecture~\ref{conj:support} asserts
that these glued structures are equinumerous with the product above; finding
the gluing bijection is the core remaining problem.

\section{Computational appendix}

All computations were performed with \texttt{qpark} (Rust, in this
repository): exact enumeration of maximal parking functions via canonical
burn orders with the skip-watermark pruning of Theorem~\ref{thm:fiber}
($Q_4$: $3\,040\,575$ maximal parking functions and $42\,467\,328$ spanning
trees in $1.5$s, matching the chromatic polynomial of $Q_4$, the Matrix-Tree
determinant, and Stanley's formula); the degree enumerator
$T_{Q_4}(1,y)$; Smith normal forms of the reduced Laplacians, giving
$K(Q_5)=\ZZ_2^5\times\ZZ_6\times\ZZ_{24}^4\times\ZZ_{48}\times
\ZZ_{192}^3\times\ZZ_{960}$ in agreement with all three theorems of Bai
\cite{bai}; and the census and spin experiments described above. Validation
is by exhaustive comparison against brute-force Dhar enumeration on small
and random multigraphs.

\begin{thebibliography}{9}
\bibitem{bai} H.~Bai, \emph{On the critical group of the $n$-cube}, Linear
Algebra Appl.\ 369 (2003) 251--261.
\bibitem{bct} B.~Benson, D.~Chakrabarty, P.~Tetali, \emph{$G$-parking
functions, acyclic orientations and spanning trees}, Discrete Math.\ 310
(2010) 1340--1353.
\bibitem{bernardi} O.~Bernardi, \emph{On the spanning trees of the hypercube
and other products of graphs}, Electron.\ J.\ Combin.\ 19(4) (2012), \#P51.
\bibitem{chari} M.~K.~Chari, \emph{Two decompositions in topological
combinatorics with applications to matroid complexes}, Trans.\ Amer.\ Math.\
Soc.\ 349 (1997) 3925--3943.
\bibitem{colmenarejo} L.~Colmenarejo, P.~E.~Harris, Z.~Jones, C.~Keller,
A.~Ramos Rodr\'iguez, E.~Sukarto, A.~R.~Vindas-Mel\'endez, \emph{Counting
$k$-Naples parking functions through permutations and the $k$-Naples area
statistic}, Enumer.\ Comb.\ Appl.\ 1(2) (2021), S2R11.
\bibitem{corilb} R.~Cori, Y.~Le~Borgne, \emph{The sand-pile model and Tutte
polynomials}, Adv.\ Appl.\ Math.\ 30 (2003) 44--52.
\bibitem{dhar} D.~Dhar, \emph{Self-organized critical state of sandpile
automaton models}, Phys.\ Rev.\ Lett.\ 64 (1990) 1613--1616.
\bibitem{ec2} R.~P.~Stanley, \emph{Enumerative Combinatorics, Volume 2},
Cambridge Univ.\ Press, 1999; 2nd ed.\ 2024, Example 5.6.10.
\bibitem{merino97} C.~Merino L\'opez, \emph{Chip firing and the Tutte
polynomial}, Ann.\ Comb.\ 1 (1997) 253--259.
\bibitem{merino01} C.~Merino, \emph{The chip firing game and matroid
complexes}, Discrete Math.\ Theor.\ Comput.\ Sci.\ Proc.\ AA (2001)
245--256.
\end{thebibliography}

\end{document}
